% ============================================================================
%  ESTADÍSTICA DESCRIPTIVA — Sesión 3
%  Semana 3: 24 al 29 de agosto de 2026
%  1.5 Tablas de distribución de frecuencia para datos no agrupados
%  1.6 Tablas de distribución de frecuencia para datos agrupados
%  Duración prevista: 3 horas
% ============================================================================
\documentclass[aspectratio=169,10pt]{beamer}

\usepackage[utf8]{inputenc}
\IfFileExists{spanish.ldf}{\usepackage[spanish,es-noshorthands]{babel}}{}
\usepackage{amsmath,amssymb}
\usepackage{booktabs}
\usepackage{tikz}
\usepackage{listings}

\usetheme{Uniminuto}

\uniasignatura{Estadística descriptiva}
\uniprofesor{Andrés Felipe Puerta Velez}
\unicorreo{andres.puerta-v@uniminuto.edu.co}
\uninrc{CBASBL031 - CBASBL141}
\unisemana{Sesión 3 — Semana 3 · 24 al 29 de agosto de 2026}

\newcommand{\listatiny}{\setbeamerfont{itemize/enumerate body}{size=\tiny}\tiny}

\newcommand{\fx}[1]{{\ttfamily\color{uniAzul}#1}}

% ------------------------------------------------------------- CÓDIGO R ----
\lstdefinestyle{Rcode}{
  language=R,
  basicstyle=\ttfamily\scriptsize,
  keywordstyle=\color{uniAzul}\bfseries,
  commentstyle=\color{uniGris}\itshape,
  stringstyle=\color{uniTitulo},
  showstringspaces=false, numbers=none, breaklines=true,
  columns=fullflexible,
  backgroundcolor=\color{uniCremaClaro},
  frame=leftline, framerule=1.2pt, rulecolor=\color{uniNaranja},
  xleftmargin=0.35cm, framexleftmargin=0.35cm,
  framexrightmargin=0.15cm, framextopmargin=2pt, framexbottommargin=2pt,
  aboveskip=0.15cm, belowskip=0.15cm
}
\lstdefinestyle{Rmini}{
  language=R,
  basicstyle=\ttfamily\tiny,
  keywordstyle=\color{uniAzul}\bfseries,
  commentstyle=\color{uniGris}\itshape,
  stringstyle=\color{uniTitulo},
  showstringspaces=false, numbers=none, breaklines=true,
  columns=fullflexible,
  backgroundcolor=\color{uniCremaClaro},
  frame=leftline, framerule=1.2pt, rulecolor=\color{uniNaranja},
  xleftmargin=0.35cm, framexleftmargin=0.35cm,
  framexrightmargin=0.15cm, framextopmargin=2pt, framexbottommargin=2pt,
  aboveskip=0.12cm, belowskip=0.12cm
}
\lstdefinestyle{Rout}{
  basicstyle=\ttfamily\tiny\color{black},
  showstringspaces=false, numbers=none, breaklines=true,
  columns=fullflexible,
  backgroundcolor=\color{uniGris!10},
  frame=leftline, framerule=1.2pt, rulecolor=\color{uniGris},
  xleftmargin=0.35cm, framexleftmargin=0.35cm,
  framexrightmargin=0.15cm, framextopmargin=2pt, framexbottommargin=2pt,
  aboveskip=0.12cm, belowskip=0.12cm
}

\begin{document}

\portada

\begin{frame}{Orden del día}
  \begin{enumerate}
    \item Llamado a lista y revisión de los ejercicios propuestos
    \item Repaso: la escala decide la herramienta
    \item \azul{Tema 1.5: tablas de frecuencia para datos no agrupados}
    \item \azul{Tema 1.6: tablas de frecuencia para datos agrupados}
    \item Taller guiado: construcción completa de una tabla
    \item Errores frecuentes y cierre
  \end{enumerate}

  \vspace{0.3cm}
  {\footnotesize\color{uniGris} Sesión de 3 horas. Trabajaremos con un mismo
  conjunto de datos durante todo el bloque y lo retomaremos en las sesiones 4 y 5.}
\end{frame}

\begin{frame}{Repaso: dónde quedamos}
  En la sesión 2 clasificamos variables y aprendimos las cuatro escalas de
  medición.

  \vspace{0.25cm}
  \begin{alertblock}{La pregunta de hoy}
    \footnotesize
    Tengo 40 datos anotados en desorden. \destacar{¿Cómo los organizo para
    poder verlos?} La respuesta es la \azul{tabla de distribución de
    frecuencias}, y la forma de construirla depende de si la variable es
    cualitativa, cuantitativa discreta o cuantitativa continua.
  \end{alertblock}

  \vspace{0.2cm}
  \begin{block}{El hilo conductor del curso}
    \footnotesize
    Datos crudos $\rightarrow$ \azul{tabla} $\rightarrow$ gráfico (sesión 4)
    $\rightarrow$ medidas resumen (unidad 2) $\rightarrow$ interpretación y
    decisión (unidad 3).
  \end{block}
\end{frame}

\begin{frame}{El caso que usaremos todo el semestre}
  \begin{exampleblock}{Almacén Bello — datos de 40 clientes}
    \footnotesize
    El administrador de un almacén en Bello quiere entender a sus clientes.
    Toma una \azul{muestra de 40} clientes del último mes y registra:

    \vspace{0.2cm}
    \begin{tabular}{@{}lll@{}}
      \toprule
      \textbf{Variable} & \textbf{Cómo se registró} & \textbf{Tipo} \\
      \midrule
      Canal de compra   & Tienda / App / Web / Teléfono & Cualitativa nominal \\
      Compras del mes   & Número entero de compras      & Cuantitativa discreta \\
      Gasto mensual     & Miles de pesos                & Cuantitativa continua \\
      \bottomrule
    \end{tabular}
  \end{exampleblock}

  \vspace{0.25cm}
  {\footnotesize\color{uniGris} Cada tipo de variable exige una tabla distinta.
  Ese es exactamente el contenido de hoy.}
\end{frame}

\seccion{1.5 Datos no agrupados}

\begin{frame}{Los cinco tipos de frecuencia}
  \scriptsize
  \begin{center}
  \begin{tabular}{@{}p{2.4cm}p{1.3cm}p{5.4cm}p{2.6cm}@{}}
    \toprule
    \textbf{Frecuencia} & \textbf{Símbolo} & \textbf{Qué es} & \textbf{Rango} \\
    \midrule
    Absoluta            & $f_i$ & Cuántas veces aparece la categoría o el valor
                        & $0 \le f_i \le n$ \\
    Absoluta acumulada  & $F_i$ & Suma de las absolutas hasta esa fila
                        & Termina en $n$ \\
    Relativa            & $h_i$ & $f_i / n$: proporción del total
                        & $0 \le h_i \le 1$ \\
    Relativa acumulada  & $H_i$ & Suma de las relativas hasta esa fila
                        & Termina en 1 \\
    Porcentual          & $\%_i$ & $h_i \times 100$
                        & Suma 100\% \\
    \bottomrule
  \end{tabular}
  \end{center}

  \vspace{0.1cm}
  \begin{alertblock}{Dos controles que nunca fallan}
    \scriptsize
    \[ \sum_{i=1}^{k} f_i = n \qquad\qquad \sum_{i=1}^{k} h_i = 1 \]
    Si su tabla no cumple estas dos igualdades, \destacar{hay un error de
    conteo}.
  \end{alertblock}
\end{frame}

\begin{frame}{Cuándo acumular y cuándo no}
  \begin{columns}[T,onlytextwidth]
    \begin{column}{0.48\textwidth}
      \begin{block}{Sí tiene sentido acumular}
        \footnotesize
        Cuando la variable tiene \azul{orden}: ordinal, discreta o continua.

        \vspace{0.15cm}
        ``El 55\% gastó menos de 232 mil'' es una frase con significado.
      \end{block}
    \end{column}
    \begin{column}{0.48\textwidth}
      \begin{alertblock}{No tiene sentido acumular}
        \footnotesize
        Cuando la variable es \azul{nominal}.

        \vspace{0.15cm}
        ``El 65\% compró por tienda o app'' depende del orden en que usted
        escribió las filas: es un número arbitrario.
      \end{alertblock}
    \end{column}
  \end{columns}

  \vspace{0.3cm}
  \begin{exampleblock}{Regla práctica}
    \footnotesize
    En una tabla de variable \destacar{nominal}, las columnas $F_i$ y $H_i$
    se omiten. Incluirlas no es un error de cálculo, pero sí de
    \azul{interpretación}, y en el examen se penaliza como tal.
  \end{exampleblock}
\end{frame}

\begin{frame}{Ejemplo guiado 1 — variable cualitativa}
  \footnotesize
  \textbf{Canal de compra de los 40 clientes} (datos crudos, en el orden en que
  se registraron):

  \vspace{0.15cm}
  \begin{beamercolorbox}[wd=\textwidth,sep=0.15cm]{block body}
    \scriptsize\ttfamily
    T A W T F A W T A T W A F T W A T A W T \\
    T W A T F A W T A T W A T W A T F A T A
  \end{beamercolorbox}
  \vspace{0.15cm}
  {\scriptsize T = Tienda física \quad A = App móvil \quad W = Web \quad
  F = Teléfono}

  \vspace{0.15cm}
  \begin{alertblock}{Paso 1: el conteo con palotes}
    \scriptsize
    Se recorre la lista \destacar{una sola vez} marcando un palote por cada
    dato. Nunca se cuenta categoría por categoría: eso obliga a recorrer la
    lista cuatro veces y multiplica los errores.
    \vspace{0.05cm}
    \begin{center}\scriptsize
    \begin{tabular}{@{}lll@{}}
      \toprule
      \textbf{Categoría} & \textbf{Conteo} & \textbf{$f_i$} \\
      \midrule
      Tienda física & \textbar\textbar\textbar\textbar\ \ \textbar\textbar\textbar\textbar\ \ \textbar\textbar\textbar\textbar & 14 \\
      App móvil     & \textbar\textbar\textbar\textbar\ \ \textbar\textbar\textbar\textbar\ \ \textbar\textbar\textbar & 13 \\
      Web           & \textbar\textbar\textbar\textbar\ \ \textbar\textbar\textbar\textbar\ & 9 \\
      Teléfono      & \textbar\textbar\textbar\textbar & 4 \\
      \bottomrule
    \end{tabular}
    \end{center}
  \end{alertblock}
\end{frame}

\begin{frame}{Ejemplo guiado 1 — la tabla terminada}
  \scriptsize
  \begin{columns}[T,onlytextwidth]
    \begin{column}{0.46\textwidth}
      \centering
      \begin{tabular}{@{}lrrr@{}}
        \toprule
        \textbf{Canal} & \textbf{$f_i$} & \textbf{$h_i$} & \textbf{$\%_i$} \\
        \midrule
        Tienda física & 14 & 0{,}350 & 35{,}0 \\
        App móvil     & 13 & 0{,}325 & 32{,}5 \\
        Web           &  9 & 0{,}225 & 22{,}5 \\
        Teléfono      &  4 & 0{,}100 & 10{,}0 \\
        \midrule
        \textbf{Total} & \textbf{40} & \textbf{1{,}000} & \textbf{100{,}0} \\
        \bottomrule
      \end{tabular}
    \end{column}
    \begin{column}{0.50\textwidth}
      \begin{exampleblock}{El cálculo, fila por fila}
        \scriptsize
        \[ h_1=\tfrac{14}{40}=0{,}350 \qquad h_2=\tfrac{13}{40}=0{,}325 \]
        \[ h_3=\tfrac{9}{40}=0{,}225 \qquad h_4=\tfrac{4}{40}=0{,}100 \]
      \end{exampleblock}
    \end{column}
  \end{columns}

  \vspace{0.15cm}
  {\scriptsize Control: $14+13+9+4=40$ \checkmark \qquad
  $0{,}350+0{,}325+0{,}225+0{,}100=1{,}000$ \checkmark}

  \vspace{0.15cm}
  \begin{alertblock}{Cómo se lee}
    \scriptsize
    \destacar{El 35\% de los clientes de la muestra compró en tienda física},
    y el canal telefónico solo representa 1 de cada 10 compras. No acumulamos
    porque el canal es \azul{nominal}.
  \end{alertblock}
\end{frame}

\begin{frame}{Ejercicio 1 — en clase}
  \begin{alertblock}{En parejas (10 minutos)}
    \footnotesize
    Una encuesta de satisfacción a 30 clientes arrojó estas respuestas:

    \vspace{0.2cm}
    \begin{beamercolorbox}[wd=0.95\textwidth,sep=0.12cm]{block body}
      \centering\scriptsize\ttfamily
      Medio Alto Medio Bajo Alto Muy alto Medio Alto Bajo Medio \\
      Alto Medio Muy bajo Alto Medio Muy alto Bajo Medio Alto Medio \\
      Muy alto Bajo Alto Muy bajo Medio Alto Muy alto Bajo Medio Muy bajo
    \end{beamercolorbox}

    \vspace{0.25cm}
    \begin{enumerate}
      \item Construya la tabla de frecuencias con $f_i$, $h_i$ y $\%_i$.
      \item ¿Esta variable admite frecuencias acumuladas? Justifique y, si la
            respuesta es sí, agregue $F_i$ y $H_i$.
      \item ¿Qué porcentaje de clientes reportó satisfacción \emph{media o
            inferior}?
    \end{enumerate}
  \end{alertblock}
\end{frame}

\begin{frame}{Ejercicio 1 — solución}
  \footnotesize
  \begin{center}
  \begin{tabular}{@{}lrrrrr@{}}
    \toprule
    \textbf{Satisfacción} & \textbf{$f_i$} & \textbf{$F_i$} & \textbf{$h_i$} &
    \textbf{$H_i$} & \textbf{$\%_i$} \\
    \midrule
    Muy bajo  &  3 &  3 & 0{,}100 & 0{,}100 & 10{,}0 \\
    Bajo      &  5 &  8 & 0{,}167 & 0{,}267 & 16{,}7 \\
    Medio     & 10 & 18 & 0{,}333 & 0{,}600 & 33{,}3 \\
    Alto      &  8 & 26 & 0{,}267 & 0{,}867 & 26{,}7 \\
    Muy alto  &  4 & 30 & 0{,}133 & 1{,}000 & 13{,}3 \\
    \midrule
    \textbf{Total} & \textbf{30} & & \textbf{1{,}000} & & \textbf{100{,}0} \\
    \bottomrule
  \end{tabular}
  \end{center}
  \vspace{0.2cm}
  \begin{exampleblock}{Puntos 2 y 3}
    \scriptsize
    \textbf{Sí} admite acumuladas: la satisfacción es \azul{ordinal}, sus
    categorías tienen un orden natural de menor a mayor. Por eso las filas
    deben escribirse en ese orden, no alfabéticamente.
    \vspace{0.1cm}
    Media o inferior: $H_3 = 0{,}600 \Rightarrow \destacar{60{,}0\%}$ de los
    clientes. Se lee directamente de la columna acumulada, sin volver a sumar.
  \end{exampleblock}
\end{frame}

\begin{frame}{Ejemplo guiado 2 — cuantitativa discreta}
  \footnotesize
  \textbf{Número de compras en el mes} de 25 clientes:

  \vspace{0.15cm}
  \begin{beamercolorbox}[wd=\textwidth,sep=0.12cm]{block body}
    \centering\scriptsize\ttfamily
    2 \ 3 \ 1 \ 0 \ 2 \ 4 \ 3 \ 2 \ 1 \ 5 \ 2 \ 3 \ 6 \ 1 \ 2 \ 4 \ 0 \ 3 \ 8 \ 2 \ 5 \ 1 \ 6 \ 4 \ 3
  \end{beamercolorbox}

  \vspace{0.25cm}
  \begin{alertblock}{Regla de decisión}
    \footnotesize
    Una variable discreta se deja \destacar{sin agrupar} cuando toma pocos
    valores distintos (guía práctica: \azul{menos de 15}). Aquí toma 8 valores
    distintos: cada uno es una fila.
  \end{alertblock}

  \vspace{0.2cm}
  {\footnotesize\color{uniGris} Ordenar los datos de menor a mayor antes de
  contar reduce los errores casi a cero. Es el primer paso siempre.}
\end{frame}

\begin{frame}{Ejemplo guiado 2 — la tabla terminada}
  \scriptsize
  \begin{columns}[T,onlytextwidth]
    \begin{column}{0.46\textwidth}
      \centering
      \begin{tabular}{@{}rrrrrr@{}}
        \toprule
        \textbf{$x_i$} & \textbf{$f_i$} & \textbf{$F_i$} & \textbf{$h_i$} &
        \textbf{$H_i$} & \textbf{$\%_i$} \\
        \midrule
        0 & 2 &  2 & 0{,}08 & 0{,}08 &  8 \\
        1 & 4 &  6 & 0{,}16 & 0{,}24 & 16 \\
        2 & 6 & 12 & 0{,}24 & 0{,}48 & 24 \\
        3 & 5 & 17 & 0{,}20 & 0{,}68 & 20 \\
        4 & 3 & 20 & 0{,}12 & 0{,}80 & 12 \\
        5 & 2 & 22 & 0{,}08 & 0{,}88 &  8 \\
        6 & 2 & 24 & 0{,}08 & 0{,}96 &  8 \\
        8 & 1 & 25 & 0{,}04 & 1{,}00 &  4 \\
        \midrule
        \textbf{Tot.} & \textbf{25} & & \textbf{1{,}00} & & \textbf{100} \\
        \bottomrule
      \end{tabular}
    \end{column}
    \begin{column}{0.50\textwidth}
      \begin{block}{Dos lecturas}
        \scriptsize
        $F_3 = 17$: \destacar{17 clientes hicieron 3 compras o menos}.

        $H_4 = 0{,}80$: el 80\% hizo 4 compras o menos.
      \end{block}
      \begin{alertblock}{La trampa}
        \scriptsize
        El valor 7 no aparece porque \azul{nadie} hizo 7 compras. Se puede
        incluir la fila con $f_i=0$ u omitirla, pero jamás se salta un valor
        que sí fue observado.
      \end{alertblock}
    \end{column}
  \end{columns}
\end{frame}

\begin{frame}{Ejercicio 2 — en clase}
  \begin{alertblock}{Individual (10 minutos)}
    \footnotesize
    Número de quejas recibidas por día durante 20 días:

    \vspace{0.2cm}
    \begin{beamercolorbox}[wd=0.9\textwidth,sep=0.12cm]{block body}
      \centering\scriptsize\ttfamily
      1 \ 3 \ 0 \ 2 \ 1 \ 4 \ 2 \ 1 \ 0 \ 3 \ 2 \ 1 \ 5 \ 2 \ 0 \ 1 \ 3 \ 2 \ 4 \ 1
    \end{beamercolorbox}

    \vspace{0.25cm}
    \begin{enumerate}
      \item Construya la tabla completa ($x_i$, $f_i$, $F_i$, $h_i$, $H_i$).
      \item ¿En cuántos días hubo \emph{a lo sumo} 2 quejas?
      \item ¿Qué porcentaje de días tuvo \emph{más} de 3 quejas?
      \item Verifique los dos controles: $\sum f_i = n$ y $\sum h_i = 1$.
    \end{enumerate}
  \end{alertblock}
\end{frame}

\begin{frame}{Ejercicio 2 — solución}
  \footnotesize
  \begin{columns}[T,onlytextwidth]
    \begin{column}{0.52\textwidth}
      \centering\scriptsize
      \begin{tabular}{@{}rrrrr@{}}
        \toprule
        \textbf{$x_i$} & \textbf{$f_i$} & \textbf{$F_i$} & \textbf{$h_i$} &
        \textbf{$H_i$} \\
        \midrule
        0 & 3 &  3 & 0{,}15 & 0{,}15 \\
        1 & 6 &  9 & 0{,}30 & 0{,}45 \\
        2 & 5 & 14 & 0{,}25 & 0{,}70 \\
        3 & 3 & 17 & 0{,}15 & 0{,}85 \\
        4 & 2 & 19 & 0{,}10 & 0{,}95 \\
        5 & 1 & 20 & 0{,}05 & 1{,}00 \\
        \midrule
        \textbf{Tot.} & \textbf{20} & & \textbf{1{,}00} & \\
        \bottomrule
      \end{tabular}
    \end{column}
    \begin{column}{0.45\textwidth}
      \begin{exampleblock}{Respuestas}
        \scriptsize
        \textbf{2.} ``A lo sumo 2'' es $F_3 = \destacar{14}$ días.

        \vspace{0.15cm}
        \textbf{3.} ``Más de 3'' es el complemento:
        $1 - H_4 = 1-0{,}85 = 0{,}15 \Rightarrow \destacar{15\%}$.

        \vspace{0.15cm}
        \textbf{4.} $3+6+5+3+2+1 = 20$ \checkmark \\
        $0{,}15+0{,}30+0{,}25+0{,}15+0{,}10+0{,}05 = 1{,}00$ \checkmark
      \end{exampleblock}
    \end{column}
  \end{columns}

  \vspace{0.15cm}
  \begin{alertblock}{El error de vocabulario más común}
    \scriptsize
    \destacar{``A lo sumo 3''} incluye el 3 ($F_i$ directo).
    \destacar{``Menos de 3''} no lo incluye (hay que usar $F_{i-1}$).
    \destacar{``Más de 3''} es $1 - H_i$. Léalo dos veces antes de responder.
  \end{alertblock}
\end{frame}

\seccion{1.6 Datos agrupados}

\begin{frame}{¿Por qué agrupar?}
  \footnotesize
  \textbf{Gasto mensual (miles de pesos) de los 40 clientes}, ya ordenado:

  \vspace{0.15cm}
  \begin{beamercolorbox}[wd=\textwidth,sep=0.12cm]{block body}
    \centering\scriptsize\ttfamily
    112 \ 127 \ 138 \ 149 \ 153 \ 158 \ 165 \ 171 \ 178 \ 184 \\
    190 \ 195 \ 199 \ 203 \ 208 \ 212 \ 215 \ 219 \ 222 \ 226 \\
    229 \ 231 \ 234 \ 238 \ 241 \ 245 \ 249 \ 253 \ 258 \ 264 \\
    269 \ 275 \ 281 \ 288 \ 295 \ 302 \ 309 \ 318 \ 331 \ 352
  \end{beamercolorbox}

  \vspace{0.25cm}
  \begin{alertblock}{El problema}
    \footnotesize
    Si hiciéramos una fila por valor tendríamos \destacar{40 filas con
    frecuencia 1} cada una: una tabla tan larga como la lista original y que no
    resume nada. Cuando la variable es continua, casi ningún valor se repite.
  \end{alertblock}

  \vspace{0.15cm}
  {\footnotesize\color{uniGris} La solución es \azul{agrupar en intervalos de
  clase}: sacrificamos un poco de detalle a cambio de poder ver la forma de la
  distribución.}
\end{frame}

\begin{frame}{El procedimiento en seis pasos}
  \footnotesize
  \begin{enumerate}\setlength{\itemsep}{0.12cm}
    \item \azul{Ordenar} los datos y localizar el mínimo y el máximo.
    \item Calcular el \azul{rango}: \quad $R = x_{\max} - x_{\min}$
    \item Decidir el \azul{número de clases} $k$ (regla de Sturges):
      \[ k = 1 + 3{,}322 \cdot \log_{10}(n) \]
    \item Calcular la \azul{amplitud} de clase: \quad $A = R / k$,
      redondeando \emph{siempre hacia arriba}.
    \item Construir los \azul{intervalos} desde $x_{\min}$, sumando $A$.
    \item \azul{Contar} cuántos datos caen en cada intervalo y completar las
      columnas de frecuencia.
  \end{enumerate}

  \vspace{0.2cm}
  \begin{alertblock}{Sobre Sturges}
    \scriptsize
    Es una \destacar{guía, no una ley}. Da un $k$ razonable para que la tabla
    no sea ni demasiado gruesa ni demasiado fina. En la práctica se acepta
    cualquier $k$ entre 5 y 15, y muchas veces se ajusta para que la amplitud
    quede en un número redondo.
  \end{alertblock}
\end{frame}

\begin{frame}{Paso a paso con el caso Almacén Bello}
  \footnotesize
  \begin{columns}[T,onlytextwidth]
    \begin{column}{0.50\textwidth}
      \begin{block}{Pasos 1 a 4}
        \scriptsize
        \textbf{1.} $x_{\min}=112$, \ $x_{\max}=352$, \ $n=40$
        \vspace{0.1cm}
        \[ \textbf{2.}\quad R = 352 - 112 = 240 \]
        \[ \textbf{3.}\quad k = 1 + 3{,}322\cdot\log_{10}(40) \]
        \[ \log_{10}(40)=1{,}6021 \]
        \[ k = 1 + 3{,}322(1{,}6021)=1+5{,}322=6{,}32 \]
        \[ \Rightarrow k = 6 \]
        \[ \textbf{4.}\quad A = \frac{R}{k} = \frac{240}{6} = 40 \]
      \end{block}
    \end{column}
    \begin{column}{0.46\textwidth}
      \begin{exampleblock}{Paso 5: los intervalos}
        \scriptsize
        Se arranca en 112 y se suma 40 seis veces:
        \vspace{0.1cm}
        \begin{tabular}{@{}ll@{}}
          Clase 1 & $[112,\ 152)$ \\
          Clase 2 & $[152,\ 192)$ \\
          Clase 3 & $[192,\ 232)$ \\
          Clase 4 & $[232,\ 272)$ \\
          Clase 5 & $[272,\ 312)$ \\
          Clase 6 & $[312,\ 352]$ \\
        \end{tabular}
      \end{exampleblock}
      \begin{alertblock}{Ojo con la última}
        \scriptsize
        Los intervalos son \destacar{cerrados por izquierda y abiertos por
        derecha}, excepto el último, que se cierra para que el máximo quede
        dentro.
      \end{alertblock}
    \end{column}
  \end{columns}
\end{frame}

\begin{frame}{Marca de clase y límites reales}
  \begin{columns}[T,onlytextwidth]
    \begin{column}{0.48\textwidth}
      \begin{block}{Marca de clase $x_i$}
        \footnotesize
        Es el punto medio del intervalo. \azul{Representa} a todos los datos de
        esa clase en los cálculos de la unidad 2.
        \[ x_i = \frac{L_{\text{inf}} + L_{\text{sup}}}{2} \]
        Ejemplo clase 3:
        \[ x_3 = \frac{192+232}{2} = 212 \]
      \end{block}
    \end{column}
    \begin{column}{0.48\textwidth}
      \begin{exampleblock}{Amplitud y verificación}
        \footnotesize
        \[ A = L_{\text{sup}} - L_{\text{inf}} = 232-192 = 40 \]
        Todas las clases deben tener la \azul{misma amplitud}. Si no, los
        gráficos y las medidas se distorsionan.

        \vspace{0.15cm}
        Verificación global:
        \[ k \cdot A = 6 \times 40 = 240 \ge R = 240 \ \checkmark \]
      \end{exampleblock}
    \end{column}
  \end{columns}

  \vspace{0.1cm}
  \begin{alertblock}{Lo que se pierde al agrupar}
    \scriptsize
    Al agrupar dejamos de saber que el cliente más gastador gastó 352. Solo
    sabemos que \destacar{alguien cayó en la clase $[312, 352]$}. Por eso los
    cálculos con datos agrupados son \azul{aproximados}.
  \end{alertblock}
\end{frame}

\begin{frame}{Paso 6: la tabla terminada}
  \scriptsize
  \begin{center}
  \begin{tabular}{@{}lrrrrrrr@{}}
    \toprule
    \textbf{Clase (gasto)} & \textbf{$x_i$} & \textbf{$f_i$} & \textbf{$F_i$} &
    \textbf{$h_i$} & \textbf{$H_i$} & \textbf{$\%_i$} & \textbf{$\%$ ac.} \\
    \midrule
    $[112,\ 152)$ & 132 &  4 &  4 & 0{,}100 & 0{,}100 & 10{,}0 &  10{,}0 \\
    $[152,\ 192)$ & 172 &  7 & 11 & 0{,}175 & 0{,}275 & 17{,}5 &  27{,}5 \\
    $[192,\ 232)$ & 212 & 11 & 22 & 0{,}275 & 0{,}550 & 27{,}5 &  55{,}0 \\
    $[232,\ 272)$ & 252 &  9 & 31 & 0{,}225 & 0{,}775 & 22{,}5 &  77{,}5 \\
    $[272,\ 312)$ & 292 &  6 & 37 & 0{,}150 & 0{,}925 & 15{,}0 &  92{,}5 \\
    $[312,\ 352]$ & 332 &  3 & 40 & 0{,}075 & 1{,}000 &  7{,}5 & 100{,}0 \\
    \midrule
    \textbf{Total} & & \textbf{40} & & \textbf{1{,}000} & & \textbf{100{,}0} & \\
    \bottomrule
  \end{tabular}
  \end{center}

  \vspace{0.1cm}
  \begin{exampleblock}{Cómo se obtuvo la fila 3, en detalle}
    \scriptsize
    Datos en $[192, 232)$: 195, 199, 203, 208, 212, 215, 219, 222, 226, 229,
    231 $\Rightarrow f_3 = 11$.
    \[ F_3 = 4+7+11 = 22 \quad h_3 = \tfrac{11}{40}=0{,}275 \quad
       H_3 = 0{,}100+0{,}175+0{,}275 = 0{,}550 \]
  \end{exampleblock}
\end{frame}

\begin{frame}{Cómo se lee esta tabla}
  \scriptsize
  \begin{block}{Cuatro frases que el administrador puede usar}
    \scriptsize
    \begin{itemize}\setlength{\itemsep}{0.04cm}
      \item La clase \azul{más frecuente} es $[192, 232)$: \destacar{27,5\%
            de los clientes} gasta entre 192 y 232 mil pesos al mes.
      \item $H_3 = 0{,}550$: el \destacar{55\%} gasta menos de 232 mil.
      \item $1 - H_4 = 1-0{,}775 = 0{,}225$: el \destacar{22,5\%} gasta 272
            mil o más. Ese es el segmento premium.
      \item $f_1 = 4$: solo 4 clientes están por debajo de 152 mil. El almacén
            casi no tiene clientes de gasto bajo.
    \end{itemize}
  \end{block}

  \vspace{0.1cm}
  \begin{alertblock}{Lo que una tabla no puede decir}
    \scriptsize
    Estas frases describen a los \azul{40 clientes de la muestra}. Afirmar que
    ``el 55\% de todos los clientes del almacén gasta menos de 232 mil''
    ya sería hablar de \destacar{clientes que nunca se midieron}. En este
    curso describimos lo observado; nos detenemos ahí.
  \end{alertblock}
\end{frame}

\begin{frame}{Ejercicio 3 — en clase (taller principal)}
  \begin{alertblock}{En grupos de 3 (20 minutos)}
    \footnotesize
    Tiempo de espera en caja, en minutos, de 30 clientes:

    \vspace{0.2cm}
    \begin{beamercolorbox}[wd=0.95\textwidth,sep=0.12cm]{block body}
      \centering\scriptsize\ttfamily
      4 \ 5 \ 7 \ 8 \ 9 \ 10 \ 11 \ 12 \ 13 \ 13 \ 14 \ 15 \ 16 \ 17 \ 17 \\
      18 \ 19 \ 20 \ 21 \ 22 \ 23 \ 24 \ 25 \ 26 \ 27 \ 28 \ 29 \ 30 \ 32 \ 34
    \end{beamercolorbox}

    \vspace{0.2cm}
    \begin{enumerate}\setlength{\itemsep}{-0.05cm}
      \item Calcule $R$, $k$ (Sturges) y $A$. Muestre el cálculo.
      \item Construya los intervalos y la tabla completa con marca de clase.
      \item ¿Qué porcentaje de clientes esperó menos de 19 minutos?
      \item ¿Qué porcentaje esperó 24 minutos o más?
      \item El gerente promete ``atención en menos de 15 minutos''.
            ¿Qué proporción de clientes recibe lo prometido?
    \end{enumerate}
  \end{alertblock}
\end{frame}

\begin{frame}{Ejercicio 3 — solución (pasos 1 y 2)}
  \footnotesize
  \begin{columns}[T,onlytextwidth]
    \begin{column}{0.42\textwidth}
      \begin{block}{Los cálculos previos}
        \scriptsize
        \[ R = 34 - 4 = 30 \]
        \[ k = 1+3{,}322\cdot\log_{10}(30) \]
        \[ \log_{10}(30) = 1{,}4771 \]
        \[ k = 1+4{,}907 = 5{,}91 \Rightarrow k=6 \]
        \[ A = \frac{30}{6} = 5 \]
      \end{block}
    \end{column}
    \begin{column}{0.56\textwidth}
      \centering\scriptsize
      \begin{tabular}{@{}lrrrrr@{}}
        \toprule
        \textbf{Clase} & \textbf{$x_i$} & \textbf{$f_i$} & \textbf{$F_i$} &
        \textbf{$h_i$} & \textbf{$H_i$} \\
        \midrule
        $[4,\ 9)$   &  6{,}5 & 4 &  4 & 0{,}133 & 0{,}133 \\
        $[9,\ 14)$  & 11{,}5 & 6 & 10 & 0{,}200 & 0{,}333 \\
        $[14,\ 19)$ & 16{,}5 & 6 & 16 & 0{,}200 & 0{,}533 \\
        $[19,\ 24)$ & 21{,}5 & 5 & 21 & 0{,}167 & 0{,}700 \\
        $[24,\ 29)$ & 26{,}5 & 5 & 26 & 0{,}167 & 0{,}867 \\
        $[29,\ 34]$ & 31{,}5 & 4 & 30 & 0{,}133 & 1{,}000 \\
        \midrule
        \textbf{Tot.} & & \textbf{30} & & \textbf{1{,}000} & \\
        \bottomrule
      \end{tabular}
    \end{column}
  \end{columns}

  \vspace{0.2cm}
  {\scriptsize\color{uniGris} Control: $4+6+6+5+5+4=30$ \checkmark \quad
  $k\cdot A = 6\times 5 = 30 \ge R$ \checkmark}
\end{frame}

\begin{frame}{Ejercicio 3 — solución (interpretación)}
  \footnotesize
  \begin{exampleblock}{Puntos 3 y 4}
    \scriptsize
    \textbf{3.} ``Menos de 19'' coincide con el límite superior de la clase 3:
    $H_3 = 0{,}533 \Rightarrow \destacar{53{,}3\%}$.

    \vspace{0.15cm}
    \textbf{4.} ``24 o más'' es el complemento de ``menos de 24'', y 24 es el
    límite superior de la clase 4:
    $1 - H_4 = 1 - 0{,}700 = 0{,}300 \Rightarrow \destacar{30{,}0\%}$.
  \end{exampleblock}

  \vspace{0.15cm}
  \begin{alertblock}{Punto 5: la pregunta difícil}
    \scriptsize
    ``Menos de 15 minutos'' \destacar{no coincide} con ningún límite de clase:
    cae dentro de $[14, 19)$. Con la tabla agrupada solo podemos
    \azul{aproximar} suponiendo que los 6 datos de esa clase se reparten de
    forma uniforme:
    \[ H(15) \approx 0{,}333 + 0{,}200 \cdot \frac{15-14}{5}
       = 0{,}333 + 0{,}040 = 0{,}373 \Rightarrow 37{,}3\% \]
    Con los datos \azul{originales} el valor exacto es $11/30 = 36{,}7\%$.
    La diferencia es \destacar{el costo de haber agrupado}.
  \end{alertblock}
\end{frame}

\begin{frame}{Errores frecuentes en el examen}
  \footnotesize
  \begin{center}
  \begin{tabular}{@{}p{5.2cm}p{6.4cm}@{}}
    \toprule
    \textbf{Error} & \textbf{Corrección} \\
    \midrule
    Intervalos que se solapan: $[112,152]$ y $[152,192]$ &
      Un dato debe caer en \azul{una sola} clase: $[112,152)$, $[152,192)$ \\
    \addlinespace
    Clases de amplitud distinta &
      Todas iguales, salvo instrucción explícita en contra \\
    \addlinespace
    Acumular frecuencias en variable nominal &
      Solo se acumula si hay orden \\
    \addlinespace
    Redondear $A$ hacia abajo &
      Siempre hacia arriba: si no, el máximo queda fuera de la tabla \\
    \addlinespace
    Reportar $h_i$ y $\%_i$ como si fueran cosas distintas &
      $\%_i = h_i \times 100$; se elige una de las dos \\
    \addlinespace
    Olvidar la fila de totales &
      Es la que permite verificar el trabajo \\
    \bottomrule
  \end{tabular}
  \end{center}
\end{frame}

\begin{frame}[fragile]{La misma tabla, en R}
  \scriptsize
  En la sesión 14 trabajaremos R a fondo. Por ahora, vea que
  \destacar{toda la tabla de hoy cabe en cinco líneas}.
\begin{lstlisting}[style=Rmini]
canal <- factor(canal)                  # cualitativa nominal
table(canal)                            # frecuencias absolutas
round(prop.table(table(canal)) * 100, 1)  # porcentajes

limites <- seq(112, 352, by = 40)       # los limites de clase
clases  <- cut(gasto, breaks = limites, right = FALSE, include.lowest = TRUE)
fi <- table(clases);  Fi <- cumsum(fi);  hi <- prop.table(fi);  Hi <- cumsum(hi)
fi
\end{lstlisting}
\begin{lstlisting}[style=Rout]
clases
[112,152) [152,192) [192,232) [232,272) [272,312) [312,352]
        4         7        11         9         6         3
\end{lstlisting}

  \vspace{-0.1cm}
  \begin{alertblock}{Los dos argumentos que importan}
    \tiny
    \fx{right = FALSE} produce intervalos \destacar{$[a,b)$}, la convención del
    curso; \fx{include.lowest = TRUE} cierra el último para que el máximo entre.
    Sin ellos R usa $(a,b]$ y la tabla no coincide con la de clase.
  \end{alertblock}
\end{frame}

\begin{frame}{Cierre de la sesión}
  \begin{block}{Lo que debe quedar claro hoy}
    \scriptsize
    \begin{itemize}\setlength{\itemsep}{0.02cm}
      \item Cinco frecuencias: $f_i$, $F_i$, $h_i$, $H_i$ y $\%_i$; y dos
            controles que siempre se verifican.
      \item Cualitativa y discreta con pocos valores: \azul{una fila por
            categoría o valor}, sin agrupar.
      \item Continua o discreta con muchos valores: \azul{intervalos de clase}
            mediante $R$, Sturges, $A$ y marca de clase.
      \item Solo se acumula cuando la variable tiene orden.
      \item Agrupar resume, pero \destacar{pierde información}: los resultados
            pasan a ser aproximados.
    \end{itemize}
  \end{block}

  \vspace{0.1cm}
  \begin{alertblock}{Trabajo independiente}
    \scriptsize
    Resolver los 20 ejercicios de la lámina siguiente y \destacar{traer la
    tabla del caso Almacén Bello impresa o en el cuaderno}: en la sesión 4 la
    convertimos en histograma, polígono y ojiva.
  \end{alertblock}
\end{frame}

\begin{frame}{Ejercicios propuestos}
  \vspace{-0.1cm}
  \begin{beamercolorbox}[wd=\textwidth,sep=0.10cm]{block body}
    \centering\scriptsize
    \textbf{1--8:} conceptos y lectura de tablas. \quad
    \textbf{9--14:} cálculo de $R$, $k$ y $A$. \quad
    \textbf{15--20:} construcción e interpretación.
  \end{beamercolorbox}
  \listatiny
  \vspace{0.15cm}
  \begin{columns}[T,onlytextwidth]
    \begin{column}{0.47\textwidth}
      \begin{enumerate}\setlength{\itemsep}{0.16cm}
        \item ¿Qué significa $F_i$?
        \item ¿Qué significa $H_i$?
        \item ¿Cuánto suma siempre $\sum h_i$?
        \item ¿Se acumula en una variable nominal? ¿Por qué?
        \item Si $n=50$ y $f_3=12$, ¿cuánto vale $h_3$?
        \item Si $H_4=0{,}62$ y $h_5=0{,}18$, ¿cuánto vale $H_5$?
        \item Si $F_5=38$ y $F_4=29$, ¿cuánto vale $f_5$?
        \item ¿Qué es la marca de clase y para qué sirve?
        \item $n=60$: calcule $k$ con Sturges.
        \item $n=100$: calcule $k$ con Sturges.
      \end{enumerate}
    \end{column}
    \begin{column}{0.50\textwidth}
      \begin{enumerate}\setlength{\itemsep}{0.16cm}
        \setcounter{enumi}{10}
        \item $x_{\min}=18$, $x_{\max}=93$: calcule $R$.
        \item Con $R=75$ y $k=7$: calcule $A$.
        \item Clase $[45, 60)$: marca de clase y amplitud.
        \item ¿Por qué $A$ se redondea hacia arriba?
        \item Con la tabla del Almacén Bello: ¿qué \% gasta menos de 192?
        \item ¿Qué \% gasta entre 192 y 272?
        \item ¿Cuántos clientes gastan 272 o más?
        \item ¿Cuál es la clase modal y qué significa?
        \item Si el gerente sube el corte premium a 312, ¿a cuántos clientes afecta?
        \item ¿Por qué agrupar hace que los cálculos sean aproximados?
      \end{enumerate}
    \end{column}
  \end{columns}
\end{frame}

\begin{frame}{Respuestas}{Autoevaluación de los ejercicios propuestos}
  \listatiny
  \vspace{0.15cm}
  \begin{columns}[T,onlytextwidth]
    \begin{column}{0.47\textwidth}
      \begin{enumerate}\setlength{\itemsep}{0.16cm}
        \item Datos con valor menor o igual al de esa fila
        \item Proporción acumulada hasta esa fila
        \item 1 (o 100\% en porcentaje)
        \item No: sin orden, la acumulada es arbitraria
        \item $12/50 = 0{,}24$
        \item $0{,}62+0{,}18 = 0{,}80$
        \item $38-29 = 9$
        \item Punto medio de la clase; representa a la clase
        \item $1+3{,}322(1{,}7782)=6{,}91 \Rightarrow 7$
        \item $1+3{,}322(2)=7{,}64 \Rightarrow 8$
      \end{enumerate}
    \end{column}
    \begin{column}{0.50\textwidth}
      \begin{enumerate}\setlength{\itemsep}{0.16cm}
        \setcounter{enumi}{10}
        \item $R = 93-18 = 75$
        \item $A = 75/7 = 10{,}71 \Rightarrow 11$
        \item $x_i = 52{,}5$; \ $A = 15$
        \item Para que $k\cdot A \ge R$ y el máximo entre en la tabla
        \item $H_2 = 27{,}5\%$
        \item $h_3+h_4 = 0{,}275+0{,}225 = 50{,}0\%$
        \item $6+3 = 9$ clientes
        \item $[192,232)$: el gasto más común de la muestra
        \item A 3 clientes ($f_6$), el 7{,}5\%
        \item Porque la marca de clase reemplaza a los datos reales
      \end{enumerate}
    \end{column}
  \end{columns}
\end{frame}

\begin{frame}{Referencias}
  \footnotesize
  \begin{itemize}
    \item[{[1]}] Monroy Saldívar, S. (2008). \textit{Estadística descriptiva}.
      Instituto Politécnico Nacional.
      \texttt{https://elibro.net/es/lc/uniminuto/titulos/74722}
    \item[{[2]}] Ordóñez Fernández, F. F. \& González Fernández, J. (2021).
      \textit{Estadística descriptiva paso a paso}. Ediciones Pirámide.
      \texttt{https://elibro.net/es/lc/uniminuto/titulos/215449}
    \item[{[3]}] Llinás Solano, H. (2014). \textit{Introducción a la
      estadística con aplicaciones en ciencias sociales}. Universidad del
      Norte. \texttt{https://elibro.net/es/lc/uniminuto/titulos/70062}
  \end{itemize}

  \vspace{0.35cm}
  \begin{center}
    {\Large\color{uniAzul}\textbf{¿Preguntas?}}
  \end{center}
\end{frame}

\end{document}
